\scanpage{16}
$A_0$ VA sur corps $k$ de car. $>0$.

Question: quand est-il possible, sur tte déf. inf. $A/\Lambda$ de $A_0$ (où $\Lambda$ $k$-algèbre artinienne locale) d'avoir une $k$-stratification sur $\operatorname{DR}^1(A/\Lambda)$, compatible avec changt de base et fonctoriel en $A$ ?

Si $M$ est la \unsure{var.} des modules formelle de $A_0$ sur $k$, \deleted{il faut} et $A_M$ la déformation universelle, il faut que $\mathcal H=\operatorname{DR}^1(A_M/M)$ ait une \uline{stratification} formelle \uline{compatible avec les opérations} $F,V$:
\diagram{\begin{tikzcd}[ampersand replacement=\&,column sep=large]
\mathcal H^{(p)} \arrow[r,shift left=.6ex,"F"] \& \mathcal H \arrow[l,shift left=.6ex,"V"]
\end{tikzcd}}
Cela implique que pour $\forall x\in M$, le type (géométrique) de $H_x$ muni de $F_x,V_x$ est le $=$ que celui de $H_0=\operatorname{DR}^1(A_0)$, muni de $F_0,V_0$. Cela implique sans doute que $A_0$ est ``ordinaire'' i.e. a le nb max. de pts d'ordre $p$, ou ce qui revient au $=$, que sur $t_{A_0^*}=H^1(A_0,\underline{\mathrm O}_{A_0})=H^{01}(A_0)$, la puissance p.ième soit un \unsure{épim.} [i.e. induise un isom de $H^{01(p)}$ sur $H^{01}$].

\margin{\figure{\begin{tikzpicture}
\node[inner sep=4pt] (d) {\begin{minipage}{9cm}
\diagram{\begin{tikzcd}[ampersand replacement=\&,column sep=small,row sep=large]
\unsure{0} \arrow[r,dash] \& H^{10} \arrow[r,dash] \arrow[d,phantom,"\deleted{\unsure{F}}" description] \& H \arrow[r,dash] \arrow[d,shift right=.5ex,"V"'] \& H^{\unsure{01}} \arrow[r] \arrow[d,phantom,"\deleted{\unsure{F}}" description] \& \unsure{0} \\
\unsure{0} \arrow[r,dash] \& H^{10(p)} \arrow[r,dash] \& H^{(p)} \arrow[u,shift right=.5ex,"F"'] \arrow[r,dash] \& H^{01(p)} \&
\end{tikzcd}}
\[
\operatorname{Im}V=\operatorname{Ker}F=H^{10(p)};\quad \operatorname{Ker}V=\operatorname{Im}F\unsure{'};
\]
\end{minipage}};
\draw (d.north east)--(d.south east)--(d.south west);
\end{tikzpicture}}}

// Mais alors \deleted{il y} pour tte $A$ sur $\Lambda$, il en est de $=$ dans $\operatorname{DR}^1(A)$. Cela implique que l'ext. $H=\operatorname{DR}^1(A)$ de $H^{01}$ par $H^{10}$ \unsure{splitte} can., par $\operatorname{Im}F=F(H^{(p)})$, et ce $\operatorname{Im}F\inserted{=H'}$ \unsure{évidemment}
\mspage{15}
